\documentclass[10pt]{article}
\usepackage[letterpaper,margin=0.8in]{geometry}
\usepackage{booktabs,array,xcolor,helvet}
\renewcommand{\familydefault}{\sfdefault}
\pagestyle{empty}
\setlength{\parindent}{0pt}\setlength{\parskip}{4pt}
\newcommand{\hl}[1]{\textbf{#1}}
\begin{document}
{\Large\bfseries WaterFire and the Fox Point Barrier: closures and the MOAs}\\[2pt]
{\small USACE New England District response to FOIA 26-326, items 3 and 4, released Oct 5, 2026. Working notes, not for publication.}

\vspace{4pt}
\hl{The two agreements produced.}
\begin{itemize}\setlength{\itemsep}{1pt}\setlength{\topsep}{2pt}
\item \hl{2014 MOA} (signed Aug 28, 2014): WaterFire pays for about 5 public events per fiscal year, estimated at \$15,000 a year, \$6,000 up front, plus any overrun and claims. Term ends Sept 30, 2018. Authority: 33 U.S.C. 701h; the Assistant Secretary of the Army (Civil Works) found the work ``not inconsistent'' on July 25, 2014.
\item \hl{2019 MOA} (dated May 31, 2019 per the file name; the date line is garbled): same terms, for 5 years from execution, so about May 2024 (inferred).
\item Not produced: records of funds received, by year (item 4 asked for them).
\end{itemize}

\hl{WaterFire gate closures, from the operations summaries.}

\vspace{2pt}
{\small
\begin{tabular}{@{}lcp{4.9in}@{}}
\toprule
Year & Closures & Agreement on file\\
\midrule
2011 & 5 & none produced\\
2012 & 7 & none produced\\
2013 & 2 & none produced\\
2014 & 6 & Jul 12 and Aug 9 come before the Aug 28 MOA (Jul 12 is also before the Jul 25 finding); 4 after\\
2015 & 6 & 2014 MOA (a flood-control closure on Aug 29 is logged ``Lower criteria to keep waterfire from flooding'')\\
2016 & 4 & 2014 MOA\\
2017 & 6 + 1 & 2014 MOA (May 27 is logged as flood control combined with a WaterFire event)\\
2018 & 4 & 2014 MOA through Sept 30 (3 closures); \hl{Nov 10 falls in a gap, Oct 1, 2018 to May 30, 2019}\\
2019 & 5 & 2019 MOA from May 31 (all 5 are after it)\\
2020--2022 & 0 & 2019 MOA (2020 and 2021 logs: ``No Waterfire Events'')\\
2023 & 1 & 2019 MOA (Jun 3; no gate times logged)\\
2024 & 1 & Jun 1, 6:36 pm to 10 pm. \hl{The 2019 MOA appears to end about a day earlier} (inferred)\\
2025 & 0 & only entry is a May 10 pump test\\
\bottomrule
\end{tabular}}

\vspace{4pt}
\hl{What to ask next.}
\begin{itemize}\setlength{\itemsep}{1pt}\setlength{\topsep}{2pt}
\item Is there an agreement or extension covering Oct 2018 to May 2019 and after May 2024?
\item What authority and payment covered the 15 WaterFire operations logged before the July 25, 2014 finding (2011--2013 and Jul 12, 2014)? Records of funds received, by year.
\item Why did WaterFire closures fall from 4--7 a year to 0 in 2020--2022?
\end{itemize}

\hl{Cautions.}
{\small Counts checked line by line against the operations summaries for every year. The 2020 and 2021 logs state ``No Waterfire Events''. 2011--2012 rows carry no gate times. The 2014 MOA was read by OCR; the MOA end dates for 2019 are inferred. Items 1, 2 and 5--8 of the request (closure criteria, design memo, drawings, protected-area map, inspection report) are still outstanding.}
\end{document}
